\chapter{\introductionname}
Las revisiones sitemáticas de la literatura (SLR) representan una herramienta esencial en la investigación académica, orientada a construir un mapa del conocimiento existente sobre un área. Una buena revisión no es una simple lista interminable de citas, sino una visión organizada de un campo de estudio. 

Su importancia radica en la capacidad para manejar las grandes cantidades de artículos que se publican constantemente sobre un tema determinado. Estas revisiones no solo permiten organizar la información, sino que además, facilitan la identificación de vacíos en la literatura, la formulación de nuevas preguntas de investigación, la evaluación de la calidad metodológica de los estudios existentes y la orientación de futuros trabajos académicos.

Las revisiones sistemáticas tambíen desempeñan un papel social dentro de la comunidad científica al ofrecer a estudiantes, investigadores y profesionales un recurso accesible para entender un tema sin tener que navegar de manera aislada a través de cientos de publicaciones, actuando en este sentido, como una guía confiable de conocimiento.

Este trabajo de fin de grado presenta una revisión exhaustiva de la
literatura proporcionando una visión actualizada y estructurada sobre el estado
del arte en dos áreas emergentes que están revolucionando el paradigma de las telecomunicaciones, la criptografía y la seguridad: las comunicaciones cuánticas y la ciberseguridad en la era post-cuántica.

Las comunicaciones cuánticas~\cite{Singh2024}, se han consolidado como un campo de
investigación de gran relevancia tanto en el ámbito académico como industrial, dado su capacidad para ofrecer canales de comunicación inherentemente seguros
basados en los principios de la física cuántica. Conceptos como la
superposición, entrelazamiento y teletransporte cuántico suponen
revolucionarias formas de transmitir y proteger la información. Eso supone un
avance importantísimo en la era actual, donde la información y los datos se han convertido en el activo más preciado de cualquier organización.

Los posibles avances del mundo cuántico plantean serios desafíos a los
sistemas criptográficos y de ciberseguridad tradicionales. Los algoritmos actuales se encuentran basados en problemas computacionales de ordenadores clásicos que podrían volverse obsoletos ante la capacidad de cómputo de los futuros ordenadores cuánticos. Este nuevo paradigma ha dado lugar al surgimiento de un área de investigación crucial, conocida como ciberseguridad en la era post-cuántica, centrada en el desarrollo de soluciones criptográficas resistentes a los ataques cuánticos~\cite{Liu2024} y la anticipación de los desafíos asociados con esta transformación tecnológica~\cite{Hussain2024}.

Esta revisión sistemática se justifica por la necesidad de proporcionar un analisis riguroso y actualizado de dos campos en constante evolución que se presentan como elementos claves para garantizar la seguridad de la informacion y resilencia de futuras infraestructuras digitales. Por ello, el estudio se centra exclusivamente en artículos científicos publicados en bases de datos académicas reconocidas. No se incluirán documentos puramente teóricos sin validación empírica ni artículos de divulgación sin respaldo metodológico. Este enfoque contribuye a una mejor comprensión del papel de las comunicaciones y la ciberseguridad en la transformación digital de las próximas décadas.

